\documentclass[12pt,letterpaper]{article}

\usepackage[T1]{fontenc}
\usepackage[margin=0.75in,headheight=1.5em]{geometry}
\usepackage{enumitem}
\usepackage{fancyhdr}
\usepackage{lastpage}
\usepackage{float}
\usepackage{tabu}
\usepackage{booktabs}
\usepackage{graphicx}
\usepackage{lmodern}
\usepackage[table]{xcolor}  
\usepackage[font=bf]{caption}

\begin{document}

\renewcommand\headrule{}

\pagestyle{fancy}
\fancyhf{}
\lfoot{COMP 3004}
\rfoot{\thepage/\pageref{LastPage}}
\cfoot{Algorithm Design Document}

% CUSTOM COMMANDS
\newcommand{\twodigits}[1]{\ifnum\value{#1}<10 0\fi\arabic{#1}}

\newcommand{\teamname}{Code First, Think Later}
\newcommand{\personone}{Kevin Hua}
\newcommand{\persontwo}{Hendrik Knoetze}
\newcommand{\personthree}{Juhandr\'e Knoetze}

%% FIGURE COMMANDS
\newcommand{\figurelabel}[1]{\label{figure:#1}}
\newcommand{\figureref}[1]{\textbf{Figure \ref{figure:#1}}}
%% END FIGURE COMMANDS

%% TABLE COMMANDS
\definecolor{thcolor}{RGB}{193,193,193}
\newcommand{\ccindent}{\hspace{1.5em}\hangindent=1.5em}
\newcommand{\tableheader}{\rowfont\bf\rowcolor{thcolor!30}}

\newcommand{\tablelabel}[1]{\label{table:#1}}
\newcommand{\tableref}[1]{\textbf{Table \ref{table:#1}}}
%% END TABLE COMMANDS

% TABLE STYLING
\everyrow{\hline}
\tabulinesep=0.5em

\setlist[itemize]{leftmargin=*,noitemsep,nolistsep}
\setlist[enumerate]{leftmargin=*,noitemsep,nolistsep}
% END TABLE STYLING

\thispagestyle{empty}

\begin{center}
	CARLETON UNIVERSITY
\end{center}

\vfill

\begin{center}
	{\fontsize{55pt}{55pt}\selectfont cuPID}
	\vspace{0.5em}\rule{\textwidth}{0.5pt}
	Algorithm Design Document
\end{center}

\vspace{5em}

\begin{center}
	\textbf{Team [\teamname{}]}\\
	\personone{}\\
	\persontwo{}\\
	\personthree{}
\end{center}

\vfill

\begin{center}
	Submitted to:\\
	Dr. Christine Laurendeau\\
	COMP 3004: Object Oriented Software Engineering\\
	School of Computer Science\\
	Carleton University
\end{center}

\vspace{2em}

\begin{center}
	\today
\end{center}

\newpage{}

\tableofcontents{}

\renewcommand{\listfigurename}{Figures}
\listoffigures

\renewcommand{\listtablename}{Tables}
\listoftables

\newpage{}

\section{Introduction}

Large-scale projects have always required a team to develop, design, and maintain. Unfortunately, it is rarely the case that team members are chosen based on their personalities or other such personal qualities - when you get right down to what's important to the employer, skills are paramount. Sure, the quirky, good-natured person might be favored over the quieter, unsocial person, but if that quieter person eclipsed the former with regards to talent, then they get chosen. Regrettably, the end result becomes a team of peerless talent that clash heavily on all other aspects. Maybe they're all `alpha' personalities and cannot handle not being either the center of attention or in charge. Maybe they despise the idea of sharing code with each other. Either way, the project suffers in quality because of this. Our firm, [Code First, Think Later], would like to propose in this document a potential algorithm that could solve all these problems.

\subsection{Purpose of Algorithm}

Our PPID (Project-Partner-IDentification) algorithm's purpose, as the name suggests, is, given a pool of potential team members, to sort those members into the optimal teams based on several factors including relevant skills, personality traits, leadership qualities, and degree of availability. The basis upon which the algorithm sorts the members is a set of qualifications given by each member. The qualifications themselves can be divided into two sections: the qualifications that they have, and the qualifications that they desire in a group member. By incorporating both these sets, we can more confidently ensure that the groups fit perfectly.

\subsection{Overview of Document}

This document will start with basic, simple details and slowly build upwards in complexity. As such, we begin with the Qualifications section where explanations and descriptions of each of our proposed qualifications are outlined. Furthermore, we have included the scale of possible responses for each of the qualifications and assigned each a unique ID to allow maximum traceability throughout this document. Following the Qualifications section, you will find the Rules section where the algorithmic logic gets explained. Once again, each rule is assigned a unique ID. Finally, the last section will pull all the parts together and explain how the rules, when applied to the qualifications, will give optimal and consistent results.

\vspace{1em}

\noindent Best regards,

\vspace{1em}

\begin{center}
	\includegraphics[scale=0.4]{imgs/logo.png} \\ \footnotesize{(Team [Code First, Think Later])}
\end{center}

\vspace{1em}

\section{Proposed Algorithm}

\subsection{Overview}

\subsection{Qualifications}

\subsubsection{Overview of Qualifications}

\begin{table}[H]
	\caption{Overview of Qualifications}\tablelabel{oq}
	\begin{tabu} to \textwidth {>{\bf}l >{\it}l X}
	    \tableheader{}Identifier & Category/Class & Description\\
		QF-01 & Skills & Grade in COMP 2401.\\
		QF-02 & Skills & Grade in COMP 2404.\\
		QF-03 & Skills & Report writing aptitude.\\
		QF-04 & Skills & Debugging aptitude.\\
		QF-05 & Leadership & Comfort in taking charge. \\
		QF-06 & Leadership & Importance of organization. \\
		QF-07 & Leadership & Decision-making ability. \\
		QF-08 & Leadership & Tendency towards initiatiation. \\
		QF-09 & Personality & Procrastination tendency. \\
		QF-10 & Personality & Sociability level. \\
		QF-11 & Personality & Preferred working style. \\
		QF-12 & Personality & Punctuality level. \\
		QF-13 & Availability & Daily availability. \\
		QF-14 & Availability & Weekly availability. \\
		QF-15 & Availability & Preferred mode of communication. \\
	\end{tabu}
\end{table}

\subsubsection{Detailed Breakdown of Qualifications}

\begin{table}[H]
	\caption{Detailed Breakdown of QF-01}\tablelabel{dq01}
	\begin{tabu} to \textwidth {>{\it}l X}
		\toprule
		Qualification Identifier & {\bf QF-01}\\
		Category & Skills \\
		Phrasing of Qualification & What grade did you receive in COMP 2401? \\
		Description & COMP 2401, Introduction to Systems Programming, is a second year course whose material covers some of the basic but essential components of systems programming. These components include, but are not limited to: memory management, program design, and data types.\\
		Scale for Responses &
		\begin{minipage}[t]{\linewidth}
			\begin{enumerate}
				\item[1.] D- to D+
				\item[2.] C- to C+
				\item[3.] B- to B+
				\item[4.] A- to A+
			\end{enumerate}
		\end{minipage}\\
		Reasoning & A strong understanding of the basic and fundamentals of systems programming is essential to creating a good, clean program. Without the skills learned in this class, software may be poorly managed with regards to memory allocations, file divisions, and use of the dreaded global variables.\\
		References & Laurendeau\cite{claurend1}\\
		\toprule
	\end{tabu}
\end{table}

\begin{table}[H]
	\caption{Detailed Breakdown of QF-02}\tablelabel{dq02}
	\begin{tabu} to \textwidth {>{\it}l X}
		\toprule
		Qualification Identifier & {\bf QF-02}\\
		Category & Skills \\
		Phrasing of Qualification & What grade did you receive in COMP 2404? \\
		Description & COMP 2404, Introduction to Software Engineering, is a crucial second year course that covers nearly all the important parts pertaining to the software engineering life-cycle. Components include: the concepts of data abstraction, UML notation, design patterns, and fault/exception handling.\\
		Scale for Responses &
		\begin{minipage}[t]{\linewidth}
			\begin{enumerate}
				\item[1.] D- to D+
				\item[2.] C- to C+
				\item[3.] B- to B+
				\item[4.] A- to A+
			\end{enumerate}
		\end{minipage}\\
		Reasoning & blah\\
		References & Laurendeau\cite{claurend2}\\
		\toprule
	\end{tabu}
\end{table}

\begin{table}[H]
	\caption{Detailed Breakdown of QF-03}\tablelabel{dq03}
	\begin{tabu} to \textwidth {>{\it}l X}
		\toprule
		Qualification Identifier & {\bf QF-03}\\
		Category & Skills \\
		Phrasing of Qualification & How confident and/or comfortable are you in report writing? \\
		Description & blah\\
		Scale for Responses &
		\begin{minipage}[t]{\linewidth}
			\begin{enumerate}
				\item[1.] Wait.. You mean writing with like, words?? 
				\item[2.] Ugh, I don't like writing.
				\item[3.] It's okay I guess.
				\item[4.] I enjoy writing!
				\item[5.] Just call me Shakespeare 2.0.
			\end{enumerate}
		\end{minipage}\\
		Reasoning & blah\\
		References & blah\\
		\toprule
	\end{tabu}
\end{table}

\begin{table}[H]
	\caption{Detailed Breakdown of QF-04}\tablelabel{dq04}
	\begin{tabu} to \textwidth {>{\it}l X}
		\toprule
		Qualification Identifier & {\bf QF-04}\\
		Category & Skills \\
		Phrasing of Qualification & How proficient are you at software debugging? \\
		Description & blah\\
		Scale for Responses &
		\begin{minipage}[t]{\linewidth}
			\begin{enumerate}
				\item[1.] You mean like flies? Yeah, I'm pretty good at killing bugs..
				\item[2.] I have no idea why it's not working right now..
				\item[3.] I can usually figure things out eventually.
				\item[4.] I have a good idea of where the bugs usually happen so I'm pretty fast at debugging!
				\item[5.] Ha. Please. I'm so good at debugging, you wouldn't even have time to notice any bugs before they're fixed.
			\end{enumerate}
		\end{minipage}\\
		Reasoning & blah\\
		References & blah\\
		\toprule
	\end{tabu}
\end{table}

\begin{table}[H]
	\caption{Detailed Breakdown of QF-05}\tablelabel{dq05}
	\begin{tabu} to \textwidth {>{\it}l X}
		\toprule
		Qualification Identifier & {\bf QF-05}\\
		Category & Leadership \\
		Phrasing of Qualification & How comfortable are in the a seat of power? \\
		Description & blah\\
		Scale for Responses &
		\begin{minipage}[t]{\linewidth}
			\begin{enumerate}
				\item[1.] Baaa. I am a sheep. Without a herder, I am helpless.
				\item[2.] I panic when people depend on me.
				\item[3.] I can handle it if I have to.
				\item[4.] I'm pretty happy taking charge of situations.
				\item[5.] DANCE FOR ME, MY PUPPETS, DANCE!\cite{himym}
			\end{enumerate}
		\end{minipage}\\
		Reasoning & blah\\
		References & blah\\
		\toprule
	\end{tabu}
\end{table}

\begin{table}[H]
	\caption{Detailed Breakdown of QF-06}\tablelabel{dq06}
	\begin{tabu} to \textwidth {>{\it}l X}
		\toprule
		Qualification Identifier & {\bf QF-06}\\
		Category & Leadership \\
		Phrasing of Qualification & Rate the level of importance you place on organization.\\
		Description & blah\\
		Scale for Responses &
		\begin{minipage}[t]{\linewidth}
			\begin{enumerate}
				\item[1.] Woah, we have a midterm today? I had no idea...
				\item[2.] I don't like putting things down because I usually can't find them again after..
				\item[3.] I have a pretty good idea of where everything is.
				\item[4.] Everything has its place and I can almost always find them there.
				\item[5.] My notes are cross-referenced and professionally labelled and filed.
			\end{enumerate}
		\end{minipage}\\
		Reasoning & blah\\
		References & blah\\
		\toprule
	\end{tabu}
\end{table}

\begin{table}[H]
	\caption{Detailed Breakdown of QF-07}\tablelabel{dq07}
	\begin{tabu} to \textwidth {>{\it}l X}
		\toprule
		Qualification Identifier & {\bf QF-07}\\
		Category & Leadership \\
		Phrasing of Qualification & Rate your comfortability in making big, important decisions. \\
		Description & blah\\
		Scale for Responses &
		\begin{minipage}[t]{\linewidth}
			\begin{enumerate}
				\item[1.] Uhh, Mom, can you help me answer this question...
				\item[2.] I am pretty indecisive.
				\item[3.] I don't really mind making decisions if I have to.
				\item[4.] I am pretty comfortable with making tough choices.
				\item[5.] I live for the thrill that you get from making enormous, life-changing decisions.
			\end{enumerate}
		\end{minipage}\\
		Reasoning & blah\\
		References & blah\\
		\toprule
	\end{tabu}
\end{table}

\begin{table}[H]
	\caption{Detailed Breakdown of QF-08}\tablelabel{dq08}
	\begin{tabu} to \textwidth {>{\it}l X}
		\toprule
		Qualification Identifier & {\bf QF-08}\\
		Category & Leadership \\
		Phrasing of Qualification & How willing are you to step up to the plate and take the initiative?\\
		Description & blah\\
		Scale for Responses &
		\begin{minipage}[t]{\linewidth}
			\begin{enumerate}
				\item[1.] Are you kidding me? What if I make a mistake.. *nail-biting intensifies*
				\item[2.] I dislike taking charge.
				\item[3.] If no one else wants to, I have no problem taking charge.
				\item[4.] I like stepping up and taking charge.
				\item[5.] Moot question - I'm always always at `the plate'.
			\end{enumerate}
		\end{minipage}\\
		Reasoning & blah\\
		References & blah\\
		\toprule
	\end{tabu}
\end{table}

\begin{table}[H]
	\caption{Detailed Breakdown of QF-09}\tablelabel{dq09}
	\begin{tabu} to \textwidth {>{\it}l X}
		\toprule
		Qualification Identifier & {\bf QF-09}\\
		Category & Personality \\
		Phrasing of Qualification & Rate your tendency to procrastinate to the last minute on work. \\
		Description & blah\\
		Scale for Responses &
		\begin{minipage}[t]{\linewidth}
			\begin{enumerate}
				\item[1.] Eh, I'll finish answering this tomorrow...
				\item[2.] If there's an entire week left, I see no reason to start now.
				\item[3.] Sometimes it happens, sometimes I start early. Depends on my mood.
				\item[4.] I prefer starting earlier rather than later.
				\item[5.] New assignment today? Please, I finished that yesterday.
			\end{enumerate}
		\end{minipage}\\
		Reasoning & blah\\
		References & blah\\
		\toprule
	\end{tabu}
\end{table}

\begin{table}[H]
	\caption{Detailed Breakdown of QF-10}\tablelabel{dq10}
	\begin{tabu} to \textwidth {>{\it}l X}
		\toprule
		Qualification Identifier & {\bf QF-10}\\
		Category & Personality \\
		Phrasing of Qualification & Rate your tendency to start/begin a discussion with relative strangers. \\
		Description & blah\\
		Scale for Responses &
		\begin{minipage}[t]{\linewidth}
			\begin{enumerate}
				\item[1.] stuff.
				\item[2.] stuff.
				\item[3.] stuff.
				\item[4.] stuff.
				\item[5.] stuff.
			\end{enumerate}
		\end{minipage}\\
		Reasoning & blah\\
		References & blah\\
		\toprule
	\end{tabu}
\end{table}

\begin{table}[H]
	\caption{Detailed Breakdown of QF-11}\tablelabel{dq11}
	\begin{tabu} to \textwidth {>{\it}l X}
		\toprule
		Qualification Identifier & {\bf QF-11}\\
		Category & Personality \\
		Phrasing of Qualification & When working on large scale projects, what is your preferred way of working? \\
		Description & blah\\
		Scale for Responses &
		\begin{minipage}[t]{\linewidth}
			\begin{enumerate}
				\item[1.] stuff.
				\item[2.] stuff.
				\item[3.] stuff.
				\item[4.] stuff.
				\item[5.] stuff.
			\end{enumerate}
		\end{minipage}\\
		Reasoning & blah\\
		References & blah\\
		\toprule
	\end{tabu}
\end{table}

\begin{table}[H]
	\caption{Detailed Breakdown of QF-12}\tablelabel{dq12}
	\begin{tabu} to \textwidth {>{\it}l X}
		\toprule
		Qualification Identifier & {\bf QF-12}\\
		Category & Personality \\
		Phrasing of Qualification & If you have a meeting tomorrow, what time would you usually arrive? \\
		Description & blah\\
		Scale for Responses &
		\begin{minipage}[t]{\linewidth}
			\begin{enumerate}
				\item[1.] stuff.
				\item[2.] stuff.
				\item[3.] stuff.
				\item[4.] stuff.
				\item[5.] stuff.
			\end{enumerate}
		\end{minipage}\\
		Reasoning & blah\\
		References & blah\\
		\toprule
	\end{tabu}
\end{table}

\begin{table}[H]
	\caption{Detailed Breakdown of QF-13}\tablelabel{dq13}
	\begin{tabu} to \textwidth {>{\it}l X}
		\toprule
		Qualification Identifier & {\bf QF-13}\\
		Category & Availability \\
		Phrasing of Qualification & During a typical day, when is your preffered working time? \\
		Description & blah\\
		Scale for Responses &
		\begin{minipage}[t]{\linewidth}
			\begin{enumerate}
				\item[1.] stuff.
				\item[2.] stuff.
				\item[3.] stuff.
				\item[4.] stuff.
				\item[5.] stuff.
			\end{enumerate}
		\end{minipage}\\
		Reasoning & blah\\
		References & blah\\
		\toprule
	\end{tabu}
\end{table}

\begin{table}[H]
	\caption{Detailed Breakdown of QF-14}\tablelabel{dq14}
	\begin{tabu} to \textwidth {>{\it}l X}
		\toprule
		Qualification Identifier & {\bf QF-14}\\
		Category & Availability \\
		Phrasing of Qualification & Based on your current schedule, which days are you usually available? \\
		Description & blah\\
		Scale for Responses &
		\begin{minipage}[t]{\linewidth}
			\begin{enumerate}
				\item[1.] stuff.
				\item[2.] stuff.
				\item[3.] stuff.
				\item[4.] stuff.
				\item[5.] stuff.
			\end{enumerate}
		\end{minipage}\\
		Reasoning & blah\\
		References & blah\\
		\toprule
	\end{tabu}
\end{table}

\begin{table}[H]
	\caption{Detailed Breakdown of QF-15}\tablelabel{dq15}
	\begin{tabu} to \textwidth {>{\it}l X}
		\toprule
		Qualification Identifier & {\bf QF-15}\\
		Category & Availability \\
		Phrasing of Qualification & Based on how often you check, what is your preferred mode of communication? \\
		Description & blah\\
		Scale for Responses &
		\begin{minipage}[t]{\linewidth}
			\begin{enumerate}
				\item[1.] stuff.
				\item[2.] stuff.
				\item[3.] stuff.
				\item[4.] stuff.
				\item[5.] stuff.
			\end{enumerate}
		\end{minipage}\\
		Reasoning & blah\\
		References & blah\\
		\toprule
	\end{tabu}
\end{table}

\subsection{Rules}
\subsubsection{Overview of Rules}

\begin{table}[H]
	\caption{Overview of Rules}\tablelabel{or}
	\begin{tabu} to \textwidth {>{\bf}l >{\it}l X}
	    \tableheader{}Identifier & Category/Class & Description\\
		R-01 &  & \\
		R-02 &  & \\
		R-03 &  & \\
		R-04 &  & \\
	\end{tabu}
\end{table}

\subsubsection{Detailed Breakdown of Rules}

\renewcommand\refname{\vskip -1cm}
\section{References}

\begin{thebibliography}{16}

\bibitem{example}
First Last, etc.
{\it Article/book Name},
Publisher, City, Year.

\bibitem{bass}
Bass, Bernard M.
{\it Leadership: Good, Better, Best},
The Free Press, 1985.

\bibitem{campion}
Campion, Michael A. et al.
{\it Relations Between Work Team Characteristics and Effectiveness: A Replication and Extension},
Personnel Psychology, West Lafayette, Indiana, 1996.

\bibitem{gachter}
Gachter, Simon et al.
{\it Who Makes a Good Leader? Cooperativeness, Optimism, and Leading-by-Example},
Economic Inquiry, 2010.

\bibitem{claurend1}
Laurendeau, Christine.
{\it COMP 2401 -- Introduction to Systems Programming: Course Outline},
Carleton University, Ottawa, Ontario, 2014.

\bibitem{morgeson}
Morgeson, Frederick P. et al.
{\it Selecting Individuals in Team Settings: The Importance of Social Skills, Personality Characteristics, and Teamwork Knowledge},
Personnel Psychology, East Lansing, Michigan, 2005.

\bibitem{podsakoff}
Podsakoff, Philip M.
{\it Transformational leader behaviors and their effects on followers' trust in leader, satisfaction, and organizational citizenship behaviors},
The Leadership Quarterly, Bloomington, Indiana, 1990.

\bibitem{rush}
Rush, Michael C. and Russel, Joyce E. A.
{\it Leader Prototypes and Prototype-Contingent Consensus in Leader Behavior Descriptions},
Journal of Experimental Social Psychology, Tennessee, 1986.

\bibitem{tjosvold}
Tjosvold, Dean et al.
{\it Conflict values and team relationships: conflict's contribution to team effectiveness and citizenship in China},
Journal of Organizational Behavior, Hong Kong, China, 2001.

\bibitem{wong}
Wong, Sze-sze and Burton, Richard M.
{\it Virtual Teams: What are their Characteristics, and Impact on Team Performance?},
Kluwer Academic Publishers, Durham, North Carolina, 2001.

\end{thebibliography}

\end{document}
